\section{Normality of the actual original projective variety}

\begin{definition}[The whole actual homogeneous quotient chart]
\label{def:actual-original-quotient-away-ring}
\lean{EnvelopeProperty.ProjSubvariety.nativeQuotientAway}
\lean{EnvelopeProperty.ProjSubvariety.quotChartNativeAwayEquiv}
\leanok
\uses{lem:pg-chart-normal,thm:fc-original-variety-quotient-native-comparison}
For the original homogeneous prime ideal $I$ and an original homogeneous
form $\sigma\notin I$ of degree $d$, the original quotient chart algebra
$A_\sigma/I_{(\sigma)}$ is isomorphic, as an entire original algebra, to
the actual homogeneous away ring of the actual grading on $S=k[x]/I$.
\end{definition}
\begin{proof}
Embed both entire rings into the same original localization $S[\bar\sigma^{-1}]$.
The original quotient-chart lift is injective because the localization of the
domain $S$ is injective and the polynomial map has exactly kernel $I$.
The actual homogeneous-away value map is injective by its quotient definition.
Both images consist precisely of homogeneous fractions
$\bar\tau/\bar\sigma^m$ with $\deg\tau=md$. For one inclusion use the proved
whole original quotient-chart image criterion. For the other, represent an
actual homogeneous-away element by a same-degree numerator and denominator,
and lift its numerator through the actual surjective original graded quotient
component map in that same degree. The localization fraction identity after
multiplication by the unit $\bar\sigma^m$ identifies the two whole images.
The two injective maps into this identical whole image give the asserted
actual algebra isomorphism. Original coefficient maps are preserved.
\end{proof}

\begin{theorem}[Actual native chart normality]
\label{thm:actual-native-projective-quotient-chart-normal}
\lean{EnvelopeProperty.ProjSubvariety.nativeQuotientAway_isIntegrallyClosed}
\leanok
\uses{def:actual-original-quotient-away-ring,lem:pg-chart-normal}
If the original homogeneous coordinate ring $S$ is integrally closed, every
same original actual native homogeneous-away ring for $\sigma\notin I$
is integrally closed.
\end{theorem}
\begin{proof}
The original projective-normality theorem makes the same original
$A_\sigma/I_{(\sigma)}$ integrally closed. Transport integral closedness
through the whole actual algebra isomorphism just constructed.
\end{proof}

\begin{theorem}[Every actual original projective local ring is normal]
\label{thm:actual-original-projective-variety-normal-stalks}
\lean{EnvelopeProperty.ProjSubvariety.nativeScheme_stalk_isIntegrallyClosed}
\lean{EnvelopeProperty.ProjSubvariety.originalNativeScheme_stalk_isIntegrallyClosed}
\lean{EnvelopeProperty.ProjSubvariety.originalNativeScheme_isGeometricallyUnibranch}
\leanok
\uses{thm:actual-native-projective-quotient-chart-normal,thm:fc-original-variety-quotient-native-comparison}
For a projectively normal original variety, every stalk of its actual quotient
Proj is integrally closed, and every stalk of the existing original closed
subscheme variety is integrally closed. The latter variety satisfies the
standard geometric-unibranch structural class.
\end{theorem}
\begin{proof}
Choose a true original affine-cover chart of the actual quotient Proj containing
the original point. Its indexing homogeneous element has an actual same-degree
polynomial lift. Membership in that chart forces the element to be nonzero,
so this polynomial lift is outside $I$. The preceding whole ring comparison
and normality theorem make this same actual affine chart ring a normal domain.
Its localization at the actual corresponding original prime is integrally
closed. Use the actual spectrum-stalk isomorphism and the actual
open-immersion stalk isomorphism to transport this fact to the same original
Proj stalk. The genuine whole quotient-Proj comparison with the existing
original closed subscheme transports it to every original variety stalk.
Those stalks are domains by the proved original integrality theorem; their
actual normality therefore derives the standard geometric-unibranch class.
No smoothness of the original variety is assumed.
\end{proof}
